\section{\sys{} Design}
\label{sec:design}
\sys{} realizes the overlap identified in \secref{sec:model} through two operations:
readiness aggregation during upload, and ordered destination notification at the data
tail. Both change placement rather than mechanism. The dependencies of
\secref{sec:epsynchronization} are preserved exactly---a consumer still waits for its
complete, visible input, and a buffer is still protected until its last reader
finishes---but the network coordination that establishes those conditions moves off the
sender's sequential control path.

Why a network element can do this and an endpoint cannot is a matter of position. An
endpoint is a consumer: it must know where bytes belong before it can accept them, so a
sender holding data with no committed destination address has nowhere to put it and
waits. A forwarding element at the fan-in point has neither constraint. It observes
every source destined for a given target before that target has committed to a layout,
and it can hold data whose final memory location is still undetermined, resolving
placement at egress. This is a property of topological position rather than of any
particular device; a programmable switch is simply the instance that adds no extra hop.

\subsection{Scope and state}
The logical \sys{} node observes all \ready{} messages and all managed data units for its synchronization domain. For each generation, it maintains a participant bitmap, sealed per-destination traffic plans, output progress, destination memory descriptors, and bounded staging. A generation identifies a layer, microbatch, communication direction, and buffer epoch. Membership is fixed for that generation; the amount of routed traffic is dynamic.

The node requires bounded data storage and a transport-aware output interface. It may terminate and re-originate a transport connection, or cooperate with a NIC that implements the ordering contract. An FPGA with attached memory is one possible implementation platform; packet-buffer capacity alone does not establish the required programmable storage or transport capability. We specify the protocol's logical operations independently of a particular wire format or ASIC pipeline.

\subsection{Resource reservation}
Before upload, the endpoint acquires a staging slot and the network reserves space for admitted early data. The slot carries its address, bounds, access rights, and epoch. Once the endpoint publishes \ready{}, it promises not to revoke or reuse those resources while the generation is live. The sender is allowed to upload into the network's reserved staging; the \sys{} node is not allowed to deliver to an unready destination.

Resource preparation can overlap prior work. If a previous consumer still owns every slot, however, the new operation waits. Multiple slots reduce this dependency while retaining explicit ownership tracking. The design trades bounded endpoint and network reservation for earlier upload in the small-batch regime.

\subsection{Readiness aggregation and early upload}
After routing and local preparation, a rank sends its readiness announcement and traffic-plan information, then immediately injects \data{}. It does not wait for peer \ready{} messages to return. The \sys{} node records each participant idempotently and stages data until all required participants are ready. It then streams staged and arriving data toward their destinations, subject to output capacity and ordering.

The landing layout is a reserved communication staging layout. Metadata that computes valid lengths and final expert offsets may run alongside transfer. If staging writes themselves require a computed layout, that computation becomes an additional release condition in the latency model.

Control and data need distinct progress guarantees: buffering the latter must not prevent the remaining \ready{} messages from opening the gate. Reserved control resources or an equivalent separation are necessary, especially if link-level backpressure is used. The node need not multicast every data packet; fan-out means routing to the planned destinations, with multicast available when supported by the data plane.

\subsection{Ordered completion notification}
For a Dispatch generation $g$, source $s$ declares its actual planned records for destination $d$. Once every relevant plan is sealed,
\begin{equation}
 N_{g,d}=\sum_s n_{g,s,d}
 \label{eq:count}
\end{equation}
is the expected output count. Zero traffic is reported explicitly. A count is expressed in the transport's actual unique records or fragment ranges, not an interchangeable mixture of token routes, multicast input packets, and output replicas.

The node appends \done$(g,d)$ when all expected records have entered the ordered output channel. This event establishes when the notification may be generated. The stronger consumption contract is
\begin{equation}
 \operatorname{observe}_d(\done(g,d))
 \Rightarrow \bigwedge_{x\in\mathcal{D}_{g,d}}\operatorname{visible}_d(x),
 \label{eq:visible}
\end{equation}
where $\mathcal{D}_{g,d}$ is the expected data set. The receiver observes the notification with the required acquire/completion semantics. A transmitted-byte total, by itself, cannot establish Equation~\eqref{eq:visible}.

Reliability and duplicate suppression must precede logical progress accounting. Otherwise, a retransmitted packet could make a counter reach $N_{g,d}$ while a different fragment is still missing. An implementation may rely on a reliable transport's unique progress or track sequence/fragment completion explicitly. If the traffic size cannot be announced in advance, an ordered end marker can seal a stream, but its cost belongs in the model.

Unlike a global barrier, \done$(g,d)$ proves only the completion of $d$'s managed input. Other destinations may still be transferring. The successor may proceed only if that local condition is sufficient. This is the central semantic condition for removing the post-exchange: replacing a stronger library synchronization requires an explicit argument that it was stronger than the application's dependency.

For Dispatch, expert computation needs its own complete input and valid layout. \done{} and the required local checks establish this condition; \recycle{} after the last reader establishes the separate overwrite condition. In the primary Direct scenario, network-generated \done{} replaces the target's dependency on the World-team notification issued after sender completion and local coordination. Source-buffer ownership and genuine global dependencies remain enforced. Equation~\eqref{eq:tailgain} measures the resulting benefit; it need not equal a fixed half RTT.

\subsection{Consumption and buffer reuse}
Dispatch receivers compact or expand valid staging records into expert input after \done{} and required local completion conditions. A count-synchronized allocation can size the final output exactly; a graph-compatible path can allocate a bounded output. Both are compatible with the same network protocol.

Combine reuses the token's contributor set established during Dispatch. A route slot, rather than a rank id alone, identifies an expert contribution, since two experts on the same rank must stay distinct. Producers upload weighted partials after local preparation. The data plane accumulates the expected fragments, returns reduced results, and appends the appropriate output completion. Input-partial completion and reduced-result output completion require separate counts. This reuses an existing class of in-network reduction mechanisms~\cite{dysharp,switchml}; its arithmetic and bandwidth benefits are not attributed to the synchronization model.

After the last local use of a staging slot, the consumer publishes \recycle{} and advances its epoch. Transport completion alone never licenses overwrite. The relevant progression is reserved, writing, visible, consuming, and reusable. An implementation can overlap distinct generations, but cannot conflate their readiness flags, plans, or output progress.

\subsection{Ordering across the network boundary}
NCCL GIN's strong-signal semantics cover preceding PUTs to the same peer on the same context~\cite{gindocs}. They do not imply completion of unrelated contexts. With multiple QPs or rails, each completion domain therefore needs its own ordered notification, followed by a NIC or endpoint join, unless a node explicitly provides a cross-domain ordering operation. That join is a real cost.

Likewise, an ordered ToR egress does not alone establish GPU visibility. A switch-generated signal must participate in a delivery contract that prevents it from becoming visible before the covered data. Reordering, retries, different receive engines, and forwarding through another GPU can otherwise invalidate the inference. If an implementation must await an additional destination acknowledgment to establish this contract, the acknowledgment's exposed delay reduces or can eliminate the proposed tail saving.

For Hybrid communication, a scale-out completion can establish arrival at a rail GPU, but not completion of subsequent NVLink forwarding. TMA completion, system-scope memory ordering, and scale-up joins remain until the managed domain explicitly covers them. A network node also cannot infer completion of a purely local expert path it never observes.

\subsection{Safety and progress obligations}
The design rests on four obligations: no endpoint write before the correct epoch is writable; no completion notification becomes observable before all covered data is visible; no slot reuse before the last consumer releases it; and no admission beyond reserved network capacity. Together with reliable delivery and duplicate suppression, these obligations ensure that a released local successor receives the complete expected input without exposing it to writes from another generation.

Staging includes both pre-release arrivals and post-release input/output imbalance. Reduction adds accumulators. When capacity is exhausted, unadmitted traffic waits at the sender and continues along the same path later. PFC can provide link-level protection~\cite{pfc}, but admission, headroom, and queue structure must preserve control progress. The protocol assumes reliable delivery; recovery from endpoint or \sys{}-node failure is outside its fault model.
